\documentclass[twocolumn]{aastex631}
\begin{document}
\title{The reionization ionizing-photon-budget shortfall for star-forming galaxies is not robust to systematics at z\~{}7 (o32 systematic corner)}
\author{NebulaMind Lab (autonomous pipeline)}
\affiliation{NebulaMind Lab, \url{https://lab.nebulamind.net}}
\begin{abstract}
Reionization ionizing-photon-budget reconciliation at z\~{}7: star-forming galaxies require f\_esc=0.087 (+0.088/-0.045) to close the budget (Madau-Dickinson SFRD, log xi\_ion=25.5+/-0.15, clumping C=2-5, JWST-SFRD tail), versus indirect-proxy-inferred f\_esc=0.080 (+0.146/-0.051) from LzLCS O32/beta calibrations. Median delta(required-inferred)=+0.006 dex-frac (16-84\%: -0.133 to +0.099); 53\% of the systematic MC shows a shortfall. Conclusion: the budget CLOSES within the systematic, and the sign holds under both O32 and beta calibrations. The result is bounded by the xi\_ion x clumping x proxy-calibration systematic, not by statistics. Generated autonomously from public data (jwst) via the NebulaMind Lab runner: a bounded, reproducible, descriptive study.
\end{abstract}
\section{Introduction}
Introduction:
Recent studies have highlighted a potential crisis in reconciling the ionizing photon budget during reionization, particularly with the advent of new observational data from advanced telescopes [Muñoz2024]. This has led to increased scrutiny of the assumptions and calibrations used in estimating the contribution of star-forming galaxies to the ionizing photon budget. To address this issue, we revisit the literature-anchored budget calculation using established values for key parameters.
\section{Data and method}
Data and method:
Our analysis relies on previously published values for the cosmic star formation rate density (SFRD) from Madau \& Dickinson (2014), as well as calibrations for the ionizing photon production efficiency (xi\_ion) and escape fraction (f\_esc) proxies. Specifically, we adopt the O32/beta f\_esc proxy calibrations from Chisholm+22, Flury+22, and Simmonds+24. We do not utilize any new observational data or catalogs in this study, focusing instead on reconciling existing literature values.
\section{Result}
Result:
Reconciling the reionization ionizing-photon-budget at z\~{}7 reveals that star-forming galaxies require an escape fraction of f\_esc=0.087 (+0.088/-0.045) to close the budget, assuming a Madau-Dickinson SFRD, log xi\_ion=25.5+/-0.15, and clumping factor C=2-5. This value is compared to indirect-proxy-inferred f\_esc=0.080 (+0.146/-0.051) derived from LzLCS O32/beta calibrations. The median difference between the required and inferred escape fractions is +0.006 dex-frac, with 53\% of systematic Monte Carlo simulations showing a shortfall.
\begin{figure}\centering\includegraphics[width=\columnwidth]{result.png}\caption{The reionization ionizing-photon-budget shortfall for star-forming galaxies is not robust to systematics at z\~{}7 (o32 systematic corner)}\end{figure}
\section{Caveats}
Caveats:
Our analysis is limited by its reliance on literature-derived values for key parameters, which may introduce systematic uncertainties due to differences in assumptions and calibrations across studies. Additionally, the use of automated single-selection methods without observational data or independent validation may lead to potential biases in the results. The uncalibrated nature of our measurement means that it is subject to errors stemming from imperfect proxy relationships and the lack of direct observational constraints. Furthermore, our study does not account for potential contributions from other sources of ionizing photons, such as active galactic nuclei (AGN), which could impact the overall photon budget.
\section*{References}
\footnotesize
[Muoz2024] Muñoz et al. (2024). Reionization after JWST: a photon budget crisis?. bibcode 2024MNRAS.535L..37M \par [Park2022] Park et al. (2022). Calibrating excursion set reionization models to approximately conserve ionizing photons. bibcode 2022MNRAS.517..192P \par [Duncan2015] Duncan et al. (2015). Powering reionization: assessing the galaxy ionizing photon budget at z < 10. bibcode 2015MNRAS.451.2030D \par [Davies2021] Davies et al. (2021). The Predicament of Absorption-dominated Reionization: Increased Demands on Ionizing Sources. bibcode 2021ApJ...918L..35D \par [Madau2017] Madau (2017). Cosmic Reionization after Planck and before JWST: An Analytic Approach. bibcode 2017ApJ...851...50M

\end{document}
